% Zumdahl Chapter 2 test -- 20 MCQ + 2 short FRQ. 45 min.
% Chapter-track assessment: covers the book's chapter, not a CED unit.
\documentclass{apchem-exam}

\apcunitword{Chapter}
\apcunit{2}
\apcunittitle{Atoms, Molecules, and Ions}
\apcdoctitle{Chapter Test}
\apcsubtitle{Zumdahl \S2.1--2.8 \textbullet\ 45 minutes \textbullet\ periodic table provided}

\begin{document}
\apctitleblock

\examsection{I}{Multiple Choice}{20 questions \textbullet\ 28 minutes \textbullet\ 1 point each}{\nocalculator}

% ---- laws / theory ---------------------------------------------------------
\begin{mcq}
  A sample of \ce{CO} contains \SI{1.00}{\gram} of C per \SI{1.33}{\gram} of
  O; a sample of \ce{CO2} contains \SI{1.00}{\gram} of C per
  \SI{2.66}{\gram} of O. These data best illustrate the law of
  \choices
    \choice{conservation of mass}
    \choice{definite proportion}
    \correct{multiple proportions}
    \choice{constant composition}
\end{mcq}
\why{Oxygen masses per fixed carbon are in a small whole-number ratio (2:1).}
\distractor{B}{that law compares samples of the SAME compound}

\begin{mcq}
  Which statement of Dalton's atomic theory required revision after the
  discovery of isotopes?
  \choices
    \correct{All atoms of a given element are identical.}
    \choice{Atoms are neither created nor destroyed in reactions.}
    \choice{Compounds contain atoms in fixed whole-number ratios.}
    \choice{Elements are composed of atoms.}
\end{mcq}
\why{Isotopes differ in neutron count and therefore mass.}

\begin{mcq}
  Lavoisier's careful weighing of sealed reaction vessels established that
  \choices
    \choice{atoms contain electrons}
    \correct{total mass is unchanged by a chemical reaction}
    \choice{compounds have fixed composition}
    \choice{atoms of one element can transmute into another}
\end{mcq}
\why{Conservation of mass.}

% ---- experiments -----------------------------------------------------------
\begin{mcq}
  In Thomson's cathode-ray experiment, the beam was deflected toward the
  positively charged plate. This shows that the particles in the beam
  \choices
    \choice{have very large mass}
    \correct{carry a negative charge}
    \choice{are emitted only by metals}
    \choice{travel faster than light}
\end{mcq}
\why{Attraction to positive means the particle is negative.}

\begin{mcq}
  Rutherford's gold-foil experiment is best described as showing that
  \choices
    \choice{all alpha particles pass straight through the foil}
    \choice{electrons are embedded in a positive sphere}
    \correct{a small fraction of alpha particles deflect at large angles}
    \choice{gold atoms are chemically inert}
\end{mcq}
\why{The rare large-angle scattering is the evidence for a tiny dense
nucleus.}
\distractor{A}{true of most particles, but that result alone fits plum
pudding too}

\begin{mcq}
  Millikan's oil-drop experiment provided the value of
  \choices
    \choice{the charge-to-mass ratio of the electron}
    \correct{the charge on a single electron}
    \choice{the mass of the proton}
    \choice{Avogadro's number}
\end{mcq}
\why{Droplet charges were multiples of one fundamental value; combined with
Thomson's $e/m$ it yields the electron's mass.}

% ---- isotope notation ------------------------------------------------------
\begin{mcq}
  How many neutrons are in \ce{^{75}_{33}As}?
  \choices
    \choice{33}
    \correct{42}
    \choice{75}
    \choice{108}
\end{mcq}
\why{$75-33 = 42$.}
\distractor{D}{added instead of subtracting}

\begin{mcq}
  The ion \ce{^{56}_{26}Fe^{3+}} contains
  \choices
    \choice{26 protons, 30 neutrons, 26 electrons}
    \correct{26 protons, 30 neutrons, 23 electrons}
    \choice{26 protons, 26 neutrons, 23 electrons}
    \choice{23 protons, 30 neutrons, 26 electrons}
\end{mcq}
\why{$3+$ means three electrons removed; the proton count never changes.}
\distractor{D}{changing protons would change the element}

\begin{mcq}
  Which species is isoelectronic with \ce{Ar}?
  \choices
    \choice{\ce{Na+}}
    \choice{\ce{F-}}
    \correct{\ce{S^2-}}
    \choice{\ce{Mg^2+}}
\end{mcq}
\why{$16+2 = 18$ electrons.}

\begin{mcq}
  An ion has 20 protons and 18 electrons. Its symbol is
  \choices
    \choice{\ce{Ar^2-}}
    \correct{\ce{Ca^2+}}
    \choice{\ce{Ca^2-}}
    \choice{\ce{K+}}
\end{mcq}
\why{20 protons is calcium; two fewer electrons than protons gives $2+$.}

% ---- ion charges -----------------------------------------------------------
\begin{mcq}
  Which pair of elements would form a compound with the formula \ce{X3Y2}?
  \choices
    \correct{a group 2A metal and a group 5A nonmetal}
    \choice{a group 1A metal and a group 6A nonmetal}
    \choice{a group 3A metal and a group 7A nonmetal}
    \choice{a group 2A metal and a group 7A nonmetal}
\end{mcq}
\why{$3(2+) + 2(3-) = 0$, e.g.\ \ce{Ca3N2}.}

\begin{mcq}
  Which metal does NOT require a Roman numeral in the names of its
  compounds?
  \choices
    \choice{iron}
    \choice{copper}
    \correct{zinc}
    \choice{lead}
\end{mcq}
\why{Zn forms only \ce{Zn^2+} (as do Ag and Al in this respect).}

% ---- nomenclature ----------------------------------------------------------
\begin{mcq}
  The correct name for \ce{Cu2S} is
  \choices
    \choice{copper sulfide}
    \choice{copper(II) sulfide}
    \correct{copper(I) sulfide}
    \choice{dicopper sulfide}
\end{mcq}
\why{\ce{S^2-} balanced by two Cu means each is \ce{Cu+}.}
\distractor{D}{prefixes are for covalent compounds only}

\begin{mcq}
  The formula for iron(III) sulfate is
  \choices
    \choice{\ce{FeSO4}}
    \choice{\ce{Fe3SO4}}
    \correct{\ce{Fe2(SO4)3}}
    \choice{\ce{Fe(SO4)3}}
\end{mcq}
\why{$2(3+) + 3(2-) = 0$.}
\distractor{B}{treats the numeral as a subscript}

\begin{mcq}
  Which name correctly matches its formula?
  \choices
    \choice{\ce{NaClO3} --- sodium chlorite}
    \correct{\ce{NaClO} --- sodium hypochlorite}
    \choice{\ce{NaClO2} --- sodium chlorate}
    \choice{\ce{NaClO4} --- sodium chlorite}
\end{mcq}
\why{The oxyanion ladder: ClO$^-$ hypochlorite, ClO$_2^-$ chlorite,
ClO$_3^-$ chlorate, ClO$_4^-$ perchlorate.}

\begin{mcq}
  The compound \ce{N2O5} is named
  \choices
    \choice{nitrogen oxide}
    \choice{nitrogen(V) oxide}
    \correct{dinitrogen pentoxide}
    \choice{nitrous oxide}
\end{mcq}
\why{Two nonmetals $\to$ prefixes.}
\distractor{B}{Roman numerals are for ionic compounds}

\begin{mcq}
  Which is the correct formula for ammonium phosphate?
  \choices
    \choice{\ce{NH4PO4}}
    \choice{\ce{(NH4)2PO4}}
    \correct{\ce{(NH4)3PO4}}
    \choice{\ce{NH4(PO4)3}}
\end{mcq}
\why{Three \ce{NH4+} balance one \ce{PO4^3-}.}

\begin{mcq}
  An aqueous acid contains the sulfite ion. Its name is
  \choices
    \choice{sulfuric acid}
    \correct{sulfurous acid}
    \choice{hydrosulfuric acid}
    \choice{persulfuric acid}
\end{mcq}
\why{\emph{-ite} anion $\to$ \emph{-ous} acid.}
\distractor{A}{that's the sulfate (\emph{-ate}) acid}

\begin{mcq}
  The name of \ce{HI(aq)} is
  \choices
    \correct{hydroiodic acid}
    \choice{iodic acid}
    \choice{iodous acid}
    \choice{hydrogen iodide acid}
\end{mcq}
\why{No oxygen in the anion $\to$ \emph{hydro-}root\emph{-ic}.}

\begin{mcq}
  Which compound is named using prefixes rather than charge balance?
  \choices
    \choice{\ce{MgBr2}}
    \choice{\ce{Fe2O3}}
    \correct{\ce{SF6}}
    \choice{\ce{K2CO3}}
\end{mcq}
\why{Two nonmetals = binary covalent.}

\examsection{II}{Free Response}{2 questions \textbullet\ 17 minutes}{\calculatorok}

% ---- FRQ 1 -----------------------------------------------------------------
\frq{short}{4}{Zumdahl \S2.2--2.4 \textbullet\ evidence and models}

Two oxides of nitrogen are analyzed. In oxide A, \SI{1.00}{\gram} of N is
combined with \SI{1.14}{\gram} of O. In oxide B, \SI{1.00}{\gram} of N is
combined with \SI{2.28}{\gram} of O.

\begin{parts}
  \item Calculate the ratio of the masses of oxygen combined with a fixed
        mass of nitrogen, and state which law the result illustrates.
        \workspace[2.2cm]
        \keyonly{\begin{apcanswer}$2.28/1.14 = 2.00$, a 2:1 ratio --- the law
        of multiple proportions.\end{apcanswer}}
  \item Explain why this result supports the existence of atoms rather than
        a continuous model of matter. \answerlines[2]{Small whole-number
        ratios arise naturally if elements combine as discrete indivisible
        units; continuous matter would allow any ratio.}
  \item If oxide A is \ce{NO}, give the formula of oxide B.
        \answerlines[1]{\ce{NO2}.}
\end{parts}

\begin{rubric}
  \rpoint{1}{(a) ratio computed as 2.00 / 2:1}
  \rpoint{1}{(a) names the law of multiple proportions}
  \rpoint{1}{(b) links whole-number ratios to discrete particles}
  \rpoint{1}{(c) \ce{NO2}}
\end{rubric}

% ---- FRQ 2 -----------------------------------------------------------------
\frq{short}{4}{Zumdahl \S2.5--2.8 \textbullet\ ions and nomenclature}

Consider the ion \ce{^{52}_{24}Cr^{3+}} and the compounds it forms.

\begin{parts}
  \item State the number of protons, neutrons, and electrons in this ion.
        \answerlines[2]{24 protons, $52-24 = 28$ neutrons, $24-3 = 21$
        electrons.}
  \item Write the formula of the compound this ion forms with the oxide ion,
        and name it. \answerlines[2]{\ce{Cr2O3}, chromium(III) oxide ---
        $2(3+) + 3(2-) = 0$.}
  \item Chromium also forms \ce{CrO}. Name it and explain why a Roman
        numeral is required for chromium but not for calcium.
        \answerlines[3]{Chromium(II) oxide. Chromium is a transition metal
        that forms more than one cation, so the numeral specifies which;
        calcium forms only \ce{Ca^2+}, so no ambiguity exists.}
\end{parts}

\begin{rubric}
  \rpoint{1}{(a) all three counts correct}
  \rpoint{1}{(b) \ce{Cr2O3} with charge-balance reasoning}
  \rpoint{1}{(b) named chromium(III) oxide}
  \rpoint{1}{(c) chromium(II) oxide AND the variable- vs fixed-charge
    explanation}
\end{rubric}

\examtotal

\end{document}
